\AppendixChapter{Dataset Schemas and Complexity-Labeling Rubric}{app:dataset}

\section{Normalized QAC Record}

\begin{lstlisting}[language={},caption={Illustrative normalized QAC record.}]
{
  "id": "stable-record-id",
  "question": "How do I complete this workflow?",
  "answer": "Reference answer",
  "context": "Grounding context",
  "complexity_label": "moderate",
  "metadata": {
    "dataset": "wixqa_expertwritten",
    "source_record_id": "source-id",
    "source_document_ids": ["article-id"],
    "label_provenance": "structural_or_outcome"
  }
}
\end{lstlisting}

Required text fields are normalized to UTF-8 and empty records are rejected.
Metadata remains extensible so source identifiers, outcome scores, judge
decisions, and split groups can be added without changing the core classifier
interface.

\section{Prepared Knowledge Document}

\begin{lstlisting}[language={},caption={Illustrative prepared WixQA document.}]
{
  "source_record_id": "article-id",
  "title": "Article title",
  "text": "Normalized article body",
  "url": "https://support.example/article",
  "article_type": "help-center",
  "source_type": "wixqa_corpus",
  "checksum": "sha256..."
}
\end{lstlisting}

Each imported row becomes one document. Chunking and embedding still occur
inside the selected Knowledge Base configuration.

\section{Labeling Rubric}

\begin{longtable}{p{0.13\textwidth}p{0.24\textwidth}p{0.49\textwidth}}
  \caption{Complexity-labeling decision rubric.}
  \label{tab:appendix-label-rubric}\\
  \toprule
  Label & Positive indicators & Exclusion and boundary guidance\\
  \midrule
  \endfirsthead
  \toprule
  Label & Positive indicators & Exclusion and boundary guidance\\
  \midrule
  \endhead
  Simple &
  General conversational request or answer available without organizational
  evidence. &
  Do not label a one-document enterprise lookup as simple merely because the
  answer is short; it is moderate when retrieval is required.\\
  Moderate &
  One article, one coherent source, one direct policy or procedural lookup. &
  Escalate to complex when evidence is split across sources or the question
  contains multiple dependent parts.\\
  Complex &
  Two or three evidence groups, multi-step synthesis, decomposition, or
  cross-passage comparison. &
  Escalate to advanced when the answer requires an iterative dependency chain,
  connector action, or recovery across tool calls.\\
  Advanced &
  Four or more hops, long conditional/exception chain, iterative research, or
  authorized external tool use. &
  Do not use advanced merely because a question is long. There must be an
  execution requirement beyond bounded multi-query aggregation.\\
  \bottomrule
\end{longtable}

\section{Outcome-Derived Silver Labels}

\begin{enumerate}
  \item Execute fixed L1, L2, L3, and L4 snapshots.
  \item Mark overlap $\geq 0.65$ as pass and overlap $<0.35$ as failure.
  \item Require faithfulness proxy $\geq 0.60$ for retrieval routes.
  \item Use the configured judge for unresolved or conflicting outcomes.
  \item Select the least-complex successful route.
  \item If unresolved, use the structural source/hop rubric.
  \item Store every metric, route outcome, judge identity, latency, cost, and
  configuration snapshot with the label.
\end{enumerate}

\section{Prepared Dataset Composition}

\begin{table}[H]
  \centering
  \caption{Selected source contribution to the four-class artifact.}
  \begin{tabular}{lrrrrr}
    \toprule
    Source & Simple & Moderate & Complex & Advanced & Total\\
    \midrule
    ExpertWritten & 0 & 148 & 52 & 0 & 200\\
    Simulated & 0 & 173 & 27 & 0 & 200\\
    Official Synthetic & 0 & 200 & 0 & 0 & 200\\
    Template supplement & 600 & 79 & 521 & 600 & 1,800\\
    \midrule
    Total & 600 & 600 & 600 & 600 & 2,400\\
    \bottomrule
  \end{tabular}
\end{table}

The source distribution shows why final reporting must separate authentic
business questions from template balancing examples. A high overall score can
hide weak generalization to ExpertWritten or Simulated records.
